\section{Local cursor semantics and staged composition}
\label{sec:cursor}
\subsection{A demand-driven abstract machine}
Let a host iterator handle $z$ support
$\langle z,\sigma\rangle\Downarrow_{\rm next}(b,z',\sigma',\alpha)$, where
$b$ is Some value or None. We count its first None as exhaustion; our client
protocol does not demand subsequent values from an exhausted cursor.
A generator poll is
$\langle g,\sigma\rangle\Downarrow_G(b,g',\sigma',\alpha)$.
The machine has Done, One($v$), Source($z$), Map($g,f$), Filter($g,p$),
Distinct($g,S$), FlatSeek($g,f$), FlatRun($g,f,z$), and Delay($e$).
Call-by-value host closures $f,p$ include their environments. The inductive
rules below define finite successful polls, including exhaustion. Absence of
such a derivation means neither a false guard nor successful exhaustion.
In particular, it does not distinguish a panic from divergence.
\begin{mathpar}

\Rule{G-Done}{}{\langle\mathsf{Done},\sigma\rangle\Downarrow_G
 (\mathsf{None},\mathsf{Done},\sigma,\emptyseq)}

\Rule{G-One}{}{\langle\mathsf{One}(v),\sigma\rangle\Downarrow_G
 (\mathsf{Some}(v),\mathsf{Done},\sigma,\emptyseq)}

\Rule{G-Source}{\langle z,\sigma\rangle\Downarrow_{\rm next}(\mathsf{Some}(v),z',\sigma',\alpha)}
 {\langle\mathsf{Source}(z),\sigma\rangle\Downarrow_G
 (\mathsf{Some}(v),\mathsf{Source}(z'),\sigma',\alpha)}

\Rule{G-Map}{\langle g,\sigma\rangle\Downarrow_G(\mathsf{Some}(v),g',\sigma_1,\alpha)\\
 \langle f(v),\sigma_1\rangle\Downarrow_H(w,\sigma_2,\beta)}
 {\langle\mathsf{Map}(g,f),\sigma\rangle\Downarrow_G
 (\mathsf{Some}(w),\mathsf{Map}(g',f),\sigma_2,\alpha\concat\beta)}

\Rule{G-Keep}{\langle g,\sigma\rangle\Downarrow_G(\mathsf{Some}(v),g',\sigma_1,\alpha)\\
 \langle p(v),\sigma_1\rangle\Downarrow_H(\mathsf{true},\sigma_2,\beta)}
 {\langle\mathsf{Filter}(g,p),\sigma\rangle\Downarrow_G
 (\mathsf{Some}(v),\mathsf{Filter}(g',p),\sigma_2,\alpha\concat\beta)}

\Rule{G-Skip}{\langle g,\sigma\rangle\Downarrow_G(\mathsf{Some}(v),g',\sigma_1,\alpha)\\
 \langle p(v),\sigma_1\rangle\Downarrow_H(\mathsf{false},\sigma_2,\beta)\quad
 \langle\mathsf{Filter}(g',p),\sigma_2\rangle\Downarrow_G(b,g'',\sigma_3,\gamma)}
 {\langle\mathsf{Filter}(g,p),\sigma\rangle\Downarrow_G
 (b,g'',\sigma_3,\alpha\concat\beta\concat\gamma)}
\end{mathpar}
Map and Filter return None when their upstream poll returns None, preserving
that poll's store and trace and becoming Done. Source likewise becomes Done
on its first None under the chosen observation protocol.
\begin{mathpar}

\Rule{G-Enter}{\langle g,\sigma\rangle\Downarrow_G(\mathsf{Some}(v),g',\sigma_1,\alpha)\\
 \langle f(v),\sigma_1\rangle\Downarrow_H(z,\sigma_2,\beta)\quad
 \langle\mathsf{FlatRun}(g',f,z),\sigma_2\rangle\Downarrow_G(b,h,\sigma_3,\gamma)}
 {\langle\mathsf{FlatSeek}(g,f),\sigma\rangle\Downarrow_G
 (b,h,\sigma_3,\alpha\concat\beta\concat\gamma)}

\Rule{G-Inner}{\langle z,\sigma\rangle\Downarrow_{\rm next}(\mathsf{Some}(v),z',\sigma',\alpha)}
 {\langle\mathsf{FlatRun}(g,f,z),\sigma\rangle\Downarrow_G
 (\mathsf{Some}(v),\mathsf{FlatRun}(g,f,z'),\sigma',\alpha)}

\Rule{G-NextOuter}{\langle z,\sigma\rangle\Downarrow_{\rm next}(\mathsf{None},z',\sigma_1,\alpha)\quad
 \langle\mathsf{FlatSeek}(g,f),\sigma_1\rangle\Downarrow_G(b,h,\sigma_2,\beta)}
 {\langle\mathsf{FlatRun}(g,f,z),\sigma\rangle\Downarrow_G(b,h,\sigma_2,\alpha\concat\beta)}

\Rule{G-Fresh}{\langle g,\sigma\rangle\Downarrow_G(\mathsf{Some}(v),g',\sigma',\alpha)\quad v\notin S}
 {\langle\mathsf{Distinct}(g,S),\sigma\rangle\Downarrow_G
 (\mathsf{Some}(v),\mathsf{Distinct}(g',S\cup\{v\}),\sigma',\alpha)}

\Rule{G-Seen}{\langle g,\sigma\rangle\Downarrow_G(\mathsf{Some}(v),g',\sigma_1,\alpha)\quad v\in S\\
 \langle\mathsf{Distinct}(g',S),\sigma_1\rangle\Downarrow_G(b,h,\sigma_2,\beta)}
 {\langle\mathsf{Distinct}(g,S),\sigma\rangle\Downarrow_G(b,h,\sigma_2,\alpha\concat\beta)}
\end{mathpar}
FlatSeek and Distinct become Done on upstream None, propagating its trace.
For $C[-]$ restricted to Map, Filter, Distinct, or FlatSeek (not FlatRun),
the terminal rules are explicitly
\begin{mathpar}
\Rule{G-Stop}{\langle g,\sigma\rangle\Downarrow_G(\mathsf{None},g',\sigma',\alpha)}
 {\langle C[g],\sigma\rangle\Downarrow_G(\mathsf{None},\mathsf{Done},\sigma',\alpha)}

\Rule{G-SourceStop}{\langle z,\sigma\rangle\Downarrow_{\rm next}(\mathsf{None},z',\sigma',\alpha)}
 {\langle\mathsf{Source}(z),\sigma\rangle\Downarrow_G(\mathsf{None},\mathsf{Done},\sigma',\alpha)}

\Rule{G-Delay}{\langle e,\sigma\rangle\Downarrow_H(g,\sigma_1,\alpha)\quad
 \langle g,\sigma_1\rangle\Downarrow_G(b,g',\sigma_2,\beta)}
 {\langle\mathsf{Delay}(e),\sigma\rangle\Downarrow_G(b,g',\sigma_2,\alpha\concat\beta)}
\end{mathpar}
These rules exhaust the finite successful cases. Distinct abstracts clone/hash/equality as
unobservable value operations; interpreting it as a Rust HashSet therefore
requires Definition~\ref{def:host}'s representation laws.

A suspended generator $\mathsf{Delay}(e)$ evaluates $e$ only on its first poll,
then polls the resulting generator in the resulting store and concatenates
traces. Construction of Delay itself evaluates no $e$. This explicitly
accounts for root-source effects and partial consumption. Dropping a cursor
performs no further generator poll; Rust destructor behavior remains part of
the host adequacy obligation.

An Access plan denotes Delay of a nested-loop cursor: first For creates a
source then maps its fresh binding; a first If uses One($()$) followed by
Filter. Each later For uses FlatSeek on the full predecessor binding with a
source closure, and each If uses Filter. The head uses Map; final Distinct
starts with an empty seen set. This is an operational definition for concrete
Rust leaves, independent of whether its logical labels are truthful.

\subsection{Access-to-Pipeline inference rules}
A lowering state $(\bar B,s,D,n)$ records ordered bindings, current stream,
emitted definitions, and next stream number. The names $s_n$ are
\code{iter}$n$. $\operatorname{tuple}(\bar B)$ denotes both the simple input
pattern and the corresponding tuple expression, according to syntactic
position. Rust leaf ASTs in the following rules are copied intact.
\begin{mathpar}

\Rule{IP-Root}{a_0=\ell\Rightarrow\mathsf{for}\ p\ \mathsf{in}\ e\quad\bar N=\binds(p)}
 {a_0\Downarrow_{\rm IP}(\bar N,s_0,[s_0=\mathsf{Scan}(p,e,(\bar N))],1)}

\Rule{IP-Seed}{\bar a=\emptyseq\ \lor\ a_0\text{ is If}}
 {\operatorname{start}(\bar a)\Downarrow_{\rm IP}(\emptyseq,s_0,[s_0=\mathsf{Unit}(())],1)}

\Rule{IP-Join}{a=\ell\Rightarrow\mathsf{for}\ q\ \mathsf{in}\ e\quad\bar N=\binds(q)}
 {(\bar B,s,D,n;a)\Rightarrow_{\rm IP}
 (\bar B\concat\bar N,s_n,D\concat[s_n=\mathsf{Join}(s,(\bar B),q,e,(\bar B\concat\bar N))],n+1)}

\Rule{IP-Filter}{a=\ell\Rightarrow\mathsf{if}\ e}
 {(\bar B,s,D,n;a)\Rightarrow_{\rm IP}
 (\bar B,s_n,D\concat[s_n=\mathsf{Filter}(s,(\bar B),e)],n+1)}

\Rule{IP-Output}{\text{all clauses consumed}}
 {(\bar B,s,D,n;e_h)\Rightarrow_{\rm IP}
 D\concat[s_n=\mathsf{Project}(s,(\bar B),e_h);
 s_{n+1}=\mathsf{Distinct}(s_n);\mathsf{return}\ s_{n+1}]}
\end{mathpar}
Root consumes the first For. Seed consumes no clause; the ordinary fold then
handles a first If. Empty typed Access values are handled by the same seed
case even though the concrete parser requires a nonempty clause list.

\begin{theorem}[Access lowering]
On Access plans with fresh simple binding patterns, IP lowering preserves
binding well-formedness and finite successful abstract poll derivations.
\end{theorem}
\begin{proof}
Root/Seed establish exactly the source binding and stream 0. Induct on the
remaining clauses. For introduces its fresh tuple names, so Join's input
pattern equals the previous binding and its output is their ordered
concatenation. If leaves that binding unchanged. Both consume the immediately
previous stream and use the next fresh stream number. Output is the unique
Project followed immediately by final Distinct, proving structural validity.
For semantics, relate an Access prefix cursor to its current pipeline stream.
Root relates the same source-expression closure and pattern map; Seed relates
the same One($()$). IP-Join introduces exactly FlatSeek with the identical
source closure and ordered binding union. IP-Filter introduces the identical
predicate closure. Thus corresponding polls use the same upstream poll and
host callback, leaving related residual cursors after Yield or exhaustion;
the relation is preserved by each induction step. Output adds identical Map
and Distinct. Delay postpones their construction on both sides. This is a
bisimulation of the finite successful poll relation, including finite
partial-demand traces with normal callback termination. It is not a theorem
about failed or divergent executions or arbitrary native Rust representations.
\end{proof}

\subsection{Pipeline-to-Rust inference rules}
The translation writes one Rust local per named definition. In this display,
method notation abbreviates the implementation's fully qualified
\code{core::iter} calls. No late emitter chooses join order.
\begin{mathpar}

\Rule{R-Unit}{}{s=\mathsf{Unit}(b)\Downarrow_R\code{let s = once(b);}}

\Rule{R-Scan}{}{s=\mathsf{Scan}(p,e,b)\Downarrow_R\\
 \code{let s = IntoIterator::into_iter(e).map(move |p| b);}}

\Rule{R-Join}{}{s=\mathsf{Join}(l,p,q,e,b)\Downarrow_R\\
 \mathsf{let}\ s=l.\mathsf{flatMap}(\mathsf{move}\ |p|\,
 \mathsf{IntoIter}(e).\mathsf{map}(\mathsf{move}\ |q|\{\operatorname{copies}(p);b\}));}

\Rule{R-Filter}{}{s=\mathsf{Filter}(l,p,e)\Downarrow_R\\
 \mathsf{let}\ s=l.\mathsf{filterMap}(\mathsf{move}\ |p|\,
 \mathsf{if}\ e\ \{\mathsf{Some}((\binds(p)))\}\ \mathsf{else}\ \{\mathsf{None}\});}

\Rule{R-Project}{}{s=\mathsf{Project}(l,p,e)\Downarrow_R
 \code{let s = l.map(move |p| e);}}

\Rule{R-Distinct}{}{s=\mathsf{Distinct}(l)\Downarrow_R\\
 \mathsf{let}\ s=\{\mathsf{let\ mut}\ seen=\mathsf{HashSet::new}();\\
 l.\mathsf{filter}(\mathsf{move}\ |row|\,seen.\mathsf{insert}(\mathsf{Clone}(row)))\};}
\end{mathpar}
$\operatorname{copies}(p)$ is, for every bound $x$ in pattern order,
\code{let x = Clone::clone(&x);}. Those copies occur once per inner item, inside
Join's inner map. Filter is implemented using consuming filter-map, not by
passing an extra borrowed tuple to the predicate. These details matter for
owned computed fields and cannot be replaced by an informal monad analogy.

\[
\Rule{R-Block}{d_j\Downarrow_R t_j\quad(0\le j\le m)}
 {\bar d;\mathsf{return}\ s_m\Downarrow_R
 \mathsf{Iterator::flatten}(\mathsf{onceWith}(\mathsf{move}\ ||\{\bar t;s_m\}))}.
\]
The outer once-with/flatten implements Delay: it constructs the named local
operators on the first poll. Standard adapter behavior is documented by
Rust~\cite{iterator}; the abstract adequacy obligation includes source and
closure ownership, pattern resolution, and Clone/HashSet laws.

\begin{theorem}[Adapter adequacy, conditional]
Assume the generated Rust type-checks; host callbacks and source iterators
realize their specified host judgments; generated cloning/hashing/dropping
are silent and representation-preserving; and clients stop after first None.
Then R-lowering simulates every finite successful abstract pipeline poll,
preserving its returned abstract values, callback trace, and exhaustion.
\end{theorem}
\begin{proof}
Relate Source to its stored native iterator, Map/Filter to their adapters,
FlatSeek/FlatRun to flat-map before/during an active inner iterator, and
Distinct's abstract set to its native HashSet. One poll of map polls upstream
once and applies its closure on Some, exactly G-Map. Consuming filter-map
polls until the first Some result of its closure or exhaustion, implementing
G-Keep/G-Skip. Flat-map's active inner
iterator implements G-Inner; exhaustion transitions to the next outer value
and constructs its source, implementing G-NextOuter/G-Enter. Prefix copies
in the inner map are silent related values by hypothesis. HashSet insert is
true exactly for unseen abstract rows, implementing G-Fresh/G-Seen; inserted
clones preserve the representation invariant. Unit, scan, and stop cases
match their source rules. The outer once-with creates the related state on
first demand and never earlier. Induct on finite poll derivations and then
on finite successful client requests to obtain the simulation.
This is a conditional adapter proof,
not a derivation of Rust's semantics or its trait laws from the proc macro.
\end{proof}

\begin{obligation}[C-Effects: complete effectful cursor semantics]
To extend these results to the full observations in Section~\ref{sec:scope},
give a small-step or coinductive semantics for divergent internal polling,
host panic/unwinding, iterator destruction, and closure-capture destruction,
and prove the corresponding simulations. Propagating a divergent host
callback alone is insufficient: filtering an infinite source with an
always-false terminating predicate diverges inside a poll without a single
divergent callback. The finite rules above make no assertion about that case.
\end{obligation}

\section{Quotation, residual contexts, and substitution of passes}
\label{sec:staging}
Write $C[\square^L_\Gamma]$ for an ordinary Rust AST context with a local hole
containing language $L$ and runtime environment $\Gamma$. This notation
belongs to the specification, not the concrete MiniLinq syntax. The enclosing
file already fixes store definitions, types, borrows, and call sites.
An Access hole eventually occupies an expression position. Typing/resolution
of the filled Rust context, not merely its isolated hole, is required.

Let $\operatorname{quote}_L(s)$ be ToTokens on an $L$ AST and
$\operatorname{parse}_L$ its exact parser. For native Rust leaves, spans and
syntax contexts are part of the representation; plain string equality is
insufficient. The applicable syntactic adequacy obligation is
$\operatorname{parse}_L(\operatorname{quote}_L(s))\equiv s$ up to the stated
normalization of delimiters/raw names while preserving resolution of leaves.
Procedural macros are not automatically hygienic~\cite{macros}.

\begin{mathpar}

\Rule{ST-Entry}{\operatorname{parse}_L(T)=s\quad s\Downarrow_f t}
 {\mathsf{entry}_{f,k}\{T\}\longrightarrow\operatorname{quote}(k(t))}

\Rule{ST-Hole}{\Gamma\vdash s:L\quad\Gamma\vdash s\Downarrow_f t:M}
 {C[\square^L_\Gamma(s)]\longrightarrow C[\square^M_\Gamma(t)]}

\Rule{ST-Rust}{\Gamma\vdash p\Downarrow_R e:\mathsf{Expr}}
 {C[\square^{\mathsf{Pipeline}}_\Gamma(p)]\longrightarrow C[e]}

\Rule{ST-Emit}{e:\mathsf{Expr}}{e\longrightarrow\operatorname{ToTokens}(e)}
\end{mathpar}
Here $k$ is a metalanguage continuation describing where the target is routed.
For current public macros it is \emph{fixed}, not a user-configurable CPS
argument. The actual continuation pairs are:
\begin{center}\small
\begin{tabular}{lll}\toprule
Entry & Typed edge & Continuation \\\midrule
\code{mini_linq} & canonical/extended CQ & \code{relational_plan} \\
\code{planned_query} & planned CQ & \code{planned_relational_plan} \\
\code{relational_plan} & base/extended index & \code{index_requirements} \\
\code{planned_relational_plan} & result indexes & \code{materialized} \\
\code{index_requirements} & basic/lazy physical & staged Rust file \\
\code{materialized} & materialized physical & staged Rust file \\
\code{pull} & Access to Pipeline & \code{iterator_pipeline} \\
\code{iterator_pipeline} & Pipeline to Expr & ToTokens \\\bottomrule
\end{tabular}
\end{center}
A user can quote a lower IR directly through its existing entry without an
upper pass. A physical pass can leave Access holes residual and lower them
later. The materialized file has one Access hole per named RA result; the
lazy/basic file has one query region. Rust leaves are escaped into those
contexts using existing quote interpolation. There is no CQ escape keyword,
new standalone wrapper, or claim that a whole shared RA DAG is one pipeline.

\subsection{The contract of a replaceable edge}
Local validity is a property of an IR term. Compilability is a property of a
term and a selected pass. Preservation is a relation between executions under
stated environment and observation contracts. These three properties must
remain separate even when a macro entry invokes adjacent contract checks.

Let $\mathsf{WF}_L$ denote locally well-formed terms of language $L$, and
let a compiler implementation $f$ be a partial computation on such terms.
Its successful result is written $f(s)=t$; an error or nontermination is not
a successful result. An edge contract specifies:
\begin{enumerate}
\item an admitted source domain $A_f\subseteq\mathsf{WF}_L$ and a
runtime precondition $\mathsf{Pre}_f$ on input instances, host stores, types,
and representations;
\item termination with a result in $\mathsf{WF}_M$ for every $s\in A_f$;
\item a representation relation between the source runtime environment and
the environment realized by the target's generated declarations;
\item a common observation projection $O$ and client class $\mathcal D$,
with $f(s)\preceq_{O,\mathcal D}s$ under that representation relation,
including the response-totality premise for the finite observation model;
\item for a residual hole, an admissible class of closing Rust contexts and
an adequacy statement permitting that refinement to be used in those contexts.
\end{enumerate}
Here $\preceq_{O,\mathcal D}$ is the behavior-set inclusion of
Section~\ref{sec:representation-observations}. Runtime preconditions can be
semantic assumptions, such as $\mathcal H_{\rm rel}$; a syntactic checker
need not decide them. The totality clause concerns compilation on the
\emph{declared} domain. It cannot be obtained by defining that domain after
the fact as whichever inputs happened to compile. A target's local IR check
also does not prove that its eventual Rust context type-checks.

Representation relations are established at the relevant runtime phase.
For example, a physical pass emits a build method whose executions must
construct stores representing the input or named relations. It does not
compute those stores during macro expansion. If a later edge consumes a
residual Access term, its assumed environment must be the same bindings,
types, borrows, and stores that this generated context actually supplies.

The contextual clause is substantive. Equality of exhausted answer sets is
not in general a congruence for a context that requests one answer and then
reads a host counter. Similarly, a lower-IR fragment that is valid with
borrowed source fields may cease to type-check when substituted in a context
that assumes owned computed fields. Composition at a language boundary thus
requires compatible observations and runtime representations, in addition to
matching AST types. This is the compiler-linking concern formalized more
generally by Patterson and Ahmed~\cite{paper-ccc}.

\begin{theorem}[Conditional composition and pass replacement]
\label{thm:composition}
Let $f:L\to M$ and $g:M\to N$ satisfy the edge contracts above for a common
$O,\mathcal D$. Suppose $f(A_f)\subseteq A_g$; the first edge establishes the
second's runtime precondition; their environment relations compose through
the same intermediate realization; and their closing contexts satisfy the
stated adequacy clauses. Then $g\circ f$ is total from $A_f$ to
$\mathsf{WF}_N$ and refines its source at $O,\mathcal D$.
If both edges establish behavior equality, the composite establishes equality.
An alternative $f'$ may replace $f$ on $A_f$ when it satisfies these same
conditions, including $f'(A_f)\subseteq A_g$.
\end{theorem}
\begin{proof}
Fix $s\in A_f$. Totality of $f$ supplies $t=f(s)\in\mathsf{WF}_M$;
domain inclusion supplies $t\in A_g$, so totality of $g$ supplies
$u=g(t)\in\mathsf{WF}_N$. Fix an admitted input, host store, and client.
The environment premise gives a single intermediate realization in which
both edge theorems apply. After the appropriate abstraction of represented
values, their observations inhabit the same space. Consequently
\[
 \{O(b):b\in\mathcal B_d(u,I,\sigma)\}
 \subseteq\{O(b):b\in\mathcal B_d(t,I,\sigma)\}
 \subseteq\{O(b):b\in\mathcal B_d(s,I,\sigma)\}.
\]
At a residual hole the contextual clauses lift these inclusions to the
admitted enclosing Rust program. Inclusion is transitive; equality is also
transitive when both inclusions are equalities. Applying the same argument
inductively proves the statement for a finite chain. The hypotheses for
$f'$ provide the same domain and observation steps.
\end{proof}

Two refinements of a source need not be equivalent to one another. A source
permitting several enumerations may admit a sorted implementation and an
implementation choosing another order. Both can satisfy a refinement
contract while an order-observing client distinguishes them. Replacement
therefore preserves the promised specification; it preserves the exact
behavior set of the old implementation only when the contract requires that
stronger equality. Conversely, a source promising sorted enumeration cannot
be lowered to arbitrary enumeration under $O_{\rm seq}$ merely because
$O_{\rm set}$ agrees.

\paragraph{Instantiation and present limits.}
For the relational route the common observation is the completed answer set
under $\mathcal H_{\rm rel}$. The local cursor edges establish a stronger
finite-poll property for their fixed operational structure. Combining these
facts yields no stronger observation than the first edge establishes.
Counterexample~\ref{cx:effects} rules out using the set theorem to justify
unrestricted effect-preserving substitution. E2 in Section~\ref{sec:gaps}
violates a different premise: a locally valid empty-binding Extend is
outside the physical lowering's actual successful domain. An implementation
contract that declares support for it has not established edge totality.

The theorem supplies a composition principle with explicit interfaces. It
does not discharge Rust contextual adequacy, synthesize an effect analysis,
or expose a configurable macro continuation. In particular, direct lower-IR
quotation makes a language boundary accessible; each quoted program must
still satisfy that boundary's own environment and stage obligations.
